\section{Lösungsvorschläge}

\subsection{Swipen im Downloadtab}
Keine der Testpersonen realisierte, dass sie das Menu der Elemente im Downloadtab mit Swipen hervorbringen können. Dies wird gelöst, indem beim erstmaligen öffnen des Downloadtabs ein Element automatisch geswiped wird. So sollten Nutzer mit dieser Funktionalität vertraut werden, ohne dass ihre Nutzung der App beeinträchtigt wird.

\subsection{Unklare Benennung}
Der \enquote{Entdecken} wird im Downloadtab zu \enquote{Verfügbar} umbenannt. Diese Beschreibung sollte für alle Nutzer verständlicher sein.

\subsection{Pin-Überlagerung}
Da der Nutzer hereinzoomen kann, um einen größeren Abstand zwischen den Pins zu bekommen wird dies nicht gelöst werden.

\subsection{Filteroptionen}
Das Hinzufügen eines Filters wird in zukünftigen Updates hinzugefügt, da es zur Zeit zu schwer zum realisieren ist.

\subsection{Hervorheben der Pins}
Ein Problem für TP4 war, dass Pins nach der Suche nicht hervorgehoben werden. Pins sollten deshalb wenn sie ausgewählt sind hervorgehoben werden, indem ihre Größe verdoppelt wird.


\subsection{Suchvorschläge}
Es war für mehrere Testpersonen problematisch, dass nicht klar war, ob die Suche funktioniert, da es kein Live Feedback gab. Dies wird gelöst, indem Suchvorschläge in beiden Suchfenstern angezeigt werden und indem das Suchmenu im Kartentab automatisch vergrößert wird, während der/die NutzerIn darin tippen. Dies ermöglicht es ihnen auch, weniger zu Tippen, weil sie den gewünschten Pin in der Liste auswählen können.